    \hspace{3mm}

\centerline{\bf \Huge  Задачи по графам.}
\centerline{\bf \Huge Уровень 1: Подлый ящер.}

\problem{1} В некотором государстве 100 городов, из которых выходит по 3 дороги, и 100 городов, из которых выходит по 4 дороги. Сколько дорог всего в этом государстве?

\problem{2} В углах доски $3\times3$ стоят кони: по 2 коня черного и белого цветов, причем кони одного цвета стоят в противоположных углах. За один ход разрешается выбрать любого коня и сделать им ход в свободную клетку. Можно ли переставить коней так, чтобы по прежнему все кони стояли в углах, но кони одного цвета стояли в углах при одной стороне?

\problem{3} Министру Гольдану необходимо составить схему, состоящую из 10 городов, соединённых между собой дорогами (одна дорога соединяет два различных города; два города может соединять не более одной дороги), при этом из одного города должно выходить 9 дорог, из одного — 8, из одного — 7, из двух — по 5 ,    из трёх — по 3 , из одного — 2, из одного — 1. Может ли министр Гольдан составить такую схему?

\problem{4} Можно ли нарисовать на плоскости 1001 отрезок так, чтобы каждый пересекался ровно с семью другими?

\problem{5} Эрдни выписал на доску первые 100 простых чисел: $2, 3, 5, 7, 11,\dotsc$ Эрдни хочет соединить эти числа рёбрами так, чтобы из каждого числа выходило столько рёбер, какова величина этого числа. Два числа могут быть соединены несколькими рёбрами. Удастся ли ему осуществить желаемое?

\problem{6} В королевстве из каждого города выходит 100 дорог, причем из каждого города можно доехать до любого другого. Одну дорогу закрыли на ремонт. Докажите, все еще от каждого города можно добраться до любого другого.

\problem{7} Докажите, что граф с $n$ вершинами, степень каждой из которых не менее $\dfrac{n-1}{2}$, связен.

\problem{8} В компании из 6 человек некоторые компаниями по трое ходили вместе в походы. Верно ли, что среди них найдутся четверо, среди которых каждые трое ходили вместе в поход, либо четверо, где никакие трое не ходили вместе в поход?

